\section*{Capítulo \ref{ch:b2:global-change} --- Mudança global e biosfera}

\begin{solution}{exo:b2:global-change:1}
$\Delta T = \lambda C_{\text{acum}}$ com $\lambda = \qty{1.65}{\celsius}$
por \qty{1000}{GtC}: \qty{1.2}{\celsius}, \qty{1.65}{\celsius},
\qty{3.3}{\celsius}.
\end{solution}

\begin{solution}{exo:b2:global-change:2}
Graus-dia: a soma, dia a dia, do excesso da temperatura média diária
sobre uma base, o \omterm{thm:b2:global-change:phenology}{tempo térmico} que governa o desenvolvimento das plantas
e dos ectotérmicos. Descompasso: quando duas espécies que interagem
regulam suas estações por sinais diferentes e o aquecimento desloca uma
e não a outra --- o papa-moscas, guiado pelo comprimento do dia na
África, chega em sua data antiga e encontra o pico de lagartas, guiado
pela temperatura, já passado há duas semanas.
\end{solution}

\begin{solution}{exo:b2:global-change:3}
Cerca de \qty{150}{m} encosta acima e \qty{150}{km} em direção ao polo
por grau. Uma espécie do cume não tem terreno mais alto: sua isoterma
sobe acima da montanha e seu hábitat desaparece, enquanto a área de cada
faixa diminui com a altitude no caminho para cima.
\end{solution}

\begin{solution}{exo:b2:global-change:4}
O \ensuremath{\mathrm{CO_2}} dissolvido forma ácido carbônico, que
libera íons hidrogênio, e o pH cai proporcionalmente ao logaritmo da
pressão de \ensuremath{\mathrm{CO_2}}. Os íons hidrogênio
adicionais convertem carbonato em bicarbonato, de modo que a
concentração de carbonato e o \omterm{thm:b2:global-change:acid}{estado de saturação} do carbonato de cálcio
caem. Os organismos que constroem aragonita --- corais, pterópodes,
muitas larvas --- e calcita --- cocolitoforídeos, foraminíferos,
moluscos --- precisam gastar mais para construir suas conchas e, abaixo
da saturação, as perdem.
\end{solution}

\begin{solution}{exo:b2:global-change:5}
\qty{1.5}{\celsius}: $1.5/1.65\times 1000 = \qty{909}{GtC}$ no total,
\qty{209}{GtC} restantes, 21 anos. \qty{2}{\celsius}:
\qty{1212}{GtC}, \qty{512}{GtC} restantes, 51 anos.
\end{solution}

\begin{solution}{exo:b2:global-change:6}
$\sqrt{2\times 150/0.12} = 50$ dias; antecipação de $1.5/0.12 = 12.5$
dias.
\end{solution}

\begin{solution}{exo:b2:global-change:7}
\qty{3}{\celsius}: as isotermas sobem $3000/6.5 = \qty{460}{m}$; a
faixa passa para \qtyrange{1260}{1860}{m}, mas a montanha termina em
\qty{1800}{m}: os \qty{60}{m} superiores da faixa se perdem e o restante
fica onde a área é $2^{-460/300} = 0.35$ do que era --- cerca de
dois terços da área perdidos. \qty{5}{\celsius}: \qty{770}{m}; a
faixa ficaria em \qtyrange{1570}{2170}{m}; só existem \qtyrange{1570}{1800}{m},
um terço da largura da faixa, a $2^{-770/300} = 0.17$ da
área --- menos de um décimo da área original: a espécie está
praticamente perdida.
\end{solution}

\begin{solution}{exo:b2:global-change:8}
560: pH $7.93$, íons hidrogênio $\times 1.9$; 800: $7.79$, $\times
2.6$; 1000: $7.70$, $\times 3.1$.
\end{solution}

\begin{solution}{exo:b2:global-change:9}
$1 - (1 - h)^{0.25}$: \qty{26}{\%}, \qty{53}{\%}, \qty{82}{\%}. A
perda imediata é menor porque o remanescente ainda abriga populações da
maioria das espécies; elas são pequenas demais para persistir e
desaparecem ao longo de décadas --- a \omterm{thm:b2:global-change:sar}{dívida de extinção}.
\end{solution}

\begin{solution}{exo:b2:global-change:10}
$c = 2\sqrt{0.7\times 20} = \qty{7.5}{km/ano}$; \qty{1000}{km} em
cerca de 130 anos. Reduzir à metade $r$ ou $D$ divide $c$ por
$\sqrt 2$, para \qty{5.3}{km/ano}; nenhum dos dois a reduz à metade ---
para reduzir a velocidade à metade é preciso dividir o produto $rD$ por quatro.
\end{solution}

\begin{solution}{exo:b2:global-change:11}
Taxa de fundo: $8\times 10^{6}\times 10^{-6} = 8$ por ano, 800 por
século; a mil vezes, 8000 por ano, \num{800000} por século. Dez mil
extinções documentadas em um século são cerca de dez vezes a taxa de
fundo, e não mil --- mas a documentação cobre apenas os poucos por cento
de espécies que estão descritas e monitoradas (aves, mamíferos, algumas
plantas), entre as quais a taxa é de fato cem vezes a de fundo ou mais;
para a maioria não descrita, a maior parte das extinções nunca é vista,
e a estimativa pela \omterm{thm:b2:global-change:sar}{relação espécie--área} é o único instrumento. A
verdade está entre o contado e o estimado, e é por isso que o intervalo
citado é tão largo.
\end{solution}

\begin{solution}{exo:b2:global-change:12}
A temperatura é proporcional às \omterm{prop:b2:global-change:tcre}{emissões acumuladas}; quando as emissões
líquidas são zero, o total acumulado para de crescer, e o aquecimento
também. No \omterm{thm:b2:biogeochemical-cycles:box}{modelo de caixa}, com emissões nulas, o excesso atmosférico
começa a cair à medida que os sumidouros continuam atuando, o que
resfriaria --- mas o oceano, ainda se aquecendo rumo ao equilíbrio com o
\ensuremath{\mathrm{CO_2}} elevado, continua a aquecer a superfície, e os
dois efeitos quase se anulam: a temperatura fica onde está durante
séculos. O que falta no argumento: os outros gases de \omterm{prop:b2:global-change:tcre}{efeito estufa} (a
curta vida do metano faz seu aquecimento cair rapidamente quando as
emissões param), os aerossóis que hoje mascaram parte do aquecimento e
desaparecem em semanas após as emissões que os produzem, e as
retroalimentações --- o carbono do permafrost, a morte das florestas ---
que poderiam acrescentar emissões próprias.
\end{solution}

\begin{solution}{pb:b2:global-change:1}
\textbf{1.} A: $700 + 80\times 10 = \qty{1500}{GtC}$. B: $700 +
\tfrac12\times 50\times 10 = \qty{950}{GtC}$.
\textbf{2.} A: $1.65\times 1.5 = \qty{2.5}{\celsius}$. B:
$1.65\times 0.95 = \qty{1.6}{\celsius}$.
\textbf{3.} 2050. A: \qty{1000}{GtC}, \qty{1.65}{\celsius}. B:
emissões $10\int_0^{30}(1 - t/50)\,\mathrm{d}t = 10(30 - 9) =
\qty{210}{GtC}$, total de 910, \qty{1.5}{\celsius}.
\textbf{4.} $\lambda\times 100 = \qty{0.165}{\celsius}$ por década,
um pouco abaixo dos $0.2$ observados (que incluem os outros gases e a
perda do mascaramento pelos aerossóis).
\textbf{5.} Sim: as \omterm{prop:b2:global-change:tcre}{emissões acumuladas} param de crescer em 2070, e o
aquecimento também, que fica em \qty{1.6}{\celsius} --- a
proporcionalidade significa que a temperatura é fixada pelo total, e não
pela taxa.
\textbf{6.} $2/1.65\times 1000 = \qty{1212}{GtC}$; a trajetória A o
atinge após $512/10 = 51$ anos, em 2071.
\textbf{7.} $\sqrt{2\times 200/0.1} = 63$ dias.
\textbf{8.} Aquecimento em relação a 2020: A $2.5 - 1.2 =
\qty{1.3}{\celsius}$, B \qty{0.4}{\celsius}. Antecipação $\Delta T/b$:
A 13 dias, B 4 dias.
\textbf{9.} Defasagem: A $20 + 13 = 33$ dias; B 24 dias.
\textbf{10.} Lagartas disponíveis do dia 15 ao dia 25 após a folheação.
B: a ave chega no dia 24, dentro da janela, por pouco. A: no dia 33,
oito dias depois da última lagarta.
\textbf{11.} A ave se antecipa $3\times 1.3 = 4$ dias: defasagem de $33 - 4 =
29$ dias --- ainda fora da janela.
\textbf{12.} A ave consegue viver na nova temperatura; o que ela não
consegue é alimentar os filhotes quando o alimento já acabou, porque seu
sinal e o de sua presa respondem ao aquecimento de maneiras diferentes.
\textbf{13.} A: $1.3/6.5\times 1000 = \qty{200}{m}$ para cima e
$1.3/0.65\times 100 = \qty{200}{km}$ em direção ao polo. B: \qty{57}{m} e
\qty{57}{km}.
\textbf{14.} Em 80 anos, a árvore se desloca \qty{8}{km}: ela não
consegue acompanhar \qty{200}{km} (A) e também fica aquém de \qty{57}{km}
(B) --- as árvores ficam para trás, e suas áreas de distribuição estarão
fora de equilíbrio com o clima durante séculos em qualquer das
trajetórias.
\textbf{15.} A: a faixa ficaria em \qtyrange{1700}{2000}{m}: ela alcança
exatamente o cume e não tem para onde subir. A montanha se estreita com a
altitude, de modo que a mesma faixa ocupa apenas $2^{-200/300} = 0.63$ vezes
a área anterior --- e um aquecimento maior não deixa terreno algum para
onde deslocá-la.
\textbf{16.} B: \qty{57}{m} para cima, fator de área $2^{-57/300} = 0.88$:
uma perda de \qty{12}{\%}.
\textbf{17.} Oito décadas a \qty{30}{km}: \qty{240}{km}, mais que os
\qty{200}{km} necessários em A: o peixe acompanha (e sai de sua antiga
zona de pesca).
\textbf{18.} Na borda inferior, a espécie encontra competidores,
patógenos e predadores vindos de baixo que antes não enfrentava, e sua
própria fisiologia, adaptada à temperatura antiga, é superada pela deles:
o recuo é tanto competitivo quanto fisiológico.
\textbf{19.} A: $\mathrm{pH} = 8.2 - 0.9\log_{10}(600/280) = 7.90$;
íons hidrogênio $10^{0.30} = 2.0$ vezes os de 1750. B: $8.04$; $1.44$
vez.
\textbf{20.} $\Omega = 4/(\text{razão})$: A $2.0$, B $2.8$. Com 2, os
corais ainda calcificam, mas lentamente, com esqueletos mais fracos e
recifes que se desgastam mais depressa do que crescem; com $2.8$ eles
constroem, embora o calor da trajetória B ainda os branqueie.
\textbf{21.} A: $400\times 0.4^{0.25} = 318$, 82 perdidas. B:
$400\times 0.8^{0.25} = 378$, 22 perdidas.
\textbf{22.} A: $318\times(1 - 0.05\times 1.3) = 298$. B:
$378\times(1 - 0.05\times 0.4) = 371$.
\textbf{23.} A: $1 - 298/400 = \qty{26}{\%}$. B: \qty{7}{\%}.
\textbf{24.} Nas duas, o desmatamento: 82 contra 20 espécies em A, 22
contra 7 em B --- o termo de aquecimento do modelo é modesto; na
realidade, os dois interagem, pois uma floresta fragmentada não consegue
deslocar sua área de distribuição.
\textbf{25.} Trajetória A: \qty{2.5}{\celsius}, pH $7.90$, um quarto das
espécies de árvores perdidas. Trajetória B: \qty{1.6}{\celsius}, pH
$8.04$, \qty{7}{\%} perdidas.
\end{solution}
